\begin{table*}[t]\centering\small
\caption{Identity of the fitted ExtraTrees comparators used in the confirmatory streams. Hashes show the first 12 hexadecimal characters; the frozen provenance JSON retains all 64. These are development-trained artifacts, distinct from the original submitted hackathon models.}
\label{tab:comparator-lineage}
\begin{tabular}{@{}llrrl@{}}\toprule
Comparator & Artifact variant & Features & Trees & Model SHA-256 prefix\\\midrule
ExtraTrees & BASELINE & 846 & 400 & \texttt{8be07d191ec9}\\
ExtraTrees short & SHORT & 1457 & 400 & \texttt{e4afa2907d5f}\\
\bottomrule\end{tabular}
\end{table*}

Both models were fitted on 32 physical parents from the
four earlier development campaigns \texttt{SIM0014--SIM0017};
validation used 8 parents and 16 observed records
from \texttt{SIM0018}. These identifiers name historical campaigns that
were already development data for this comparator construction.
The fitted label set contains normal operation and all 18 catalog
sources. The training target is the active source; after restoration,
including physical recovery, its label is normal. This is not a
source-excluded transfer test.

The confirmatory loader selects the two \texttt{M03\_R09} artifacts
in Table~\ref{tab:comparator-lineage}, not the original hackathon
serializations. Each report selects the largest tree vote, and six
consecutive reports of the same label drive the episode interface.
The wrapper has no separate probability-trigger threshold to lower.
Tree votes are uncalibrated; the 300-report prefix and monitoring
start at zero-based report 360 are frozen.

Full model and source hashes, training metadata and the observed
class set are recorded in \nolinkurl{generated/comparator_lineage.json}.
The reproduction snapshot is
\nolinkurl{generated/frozen/comparator_lineage_20261103/}.
Running \nolinkurl{scripts/build_comparator_lineage.py} regenerates
this table from that snapshot without the research checkout, model
fitting, predictions or campaign outputs.
